\documentclass[10pt,a4paper]{amsart}
\usepackage{amsmath,amssymb,amsthm,graphicx}
\usepackage{verbatim}
\newenvironment{reference}{\comment}{\endcomment}
%\newenvironment{reference}{}{}
\newenvironment{slogan}{\comment}{\endcomment}
\newenvironment{history}{\comment}{\endcomment}
% For commutative diagrams we use Xy-pic
\usepackage[all]{xy}

% We use 2cell for 2-commutative diagrams.
\xyoption{2cell}
% Theorem environments.
%
\theoremstyle{plain}
\newtheorem{theorem}[subsection]{Theorem}
\newtheorem{proposition}[subsection]{Proposition}
\newtheorem{lemma}[subsection]{Lemma}

\theoremstyle{definition}
\newtheorem{definition}[subsection]{Definition}
\newtheorem{example}[subsection]{Example}
\newtheorem{exercise}[subsection]{Exercise}
\newtheorem{situation}[subsection]{Situation}

\theoremstyle{remark}
\newtheorem{remark}[subsection]{Remark}
\newtheorem{remarks}[subsection]{Remarks}

\numberwithin{equation}{subsection}

% Macros
%
\def\lim{\mathop{\mathrm{lim}}\nolimits}
\def\colim{\mathop{\mathrm{colim}}\nolimits}
\def\Spec{\mathop{\mathrm{Spec}}}
\def\Hom{\mathop{\mathrm{Hom}}\nolimits}
\def\Ext{\mathop{\mathrm{Ext}}\nolimits}
\def\SheafHom{\mathop{\mathcal{H}\!\mathit{om}}\nolimits}
\def\SheafExt{\mathop{\mathcal{E}\!\mathit{xt}}\nolimits}
\def\Sch{\mathit{Sch}}
\def\Mor{\mathop{\mathrm{Mor}}\nolimits}
\def\Ob{\mathop{\mathrm{Ob}}\nolimits}
\def\Sh{\mathop{\mathit{Sh}}\nolimits}
\def\NL{\mathop{N\!L}\nolimits}
\def\CH{\mathop{\mathrm{CH}}\nolimits}
\def\proetale{{pro\text{-}\acute{e}tale}}
\def\etale{{\acute{e}tale}}
\def\QCoh{\mathit{QCoh}}
\def\Ker{\mathop{\mathrm{Ker}}}
\def\Im{\mathop{\mathrm{Im}}}
\def\Coker{\mathop{\mathrm{Coker}}}
\def\Coim{\mathop{\mathrm{Coim}}}

% Boxtimes
%
\DeclareMathSymbol{\boxtimes}{\mathbin}{AMSa}{"02}

%
% Macros for moduli stacks/spaces
%
\def\QCohstack{\mathcal{QC}\!\mathit{oh}}
\def\Cohstack{\mathcal{C}\!\mathit{oh}}
\def\Spacesstack{\mathcal{S}\!\mathit{paces}}
\def\Quotfunctor{\mathrm{Quot}}
\def\Hilbfunctor{\mathrm{Hilb}}
\def\Curvesstack{\mathcal{C}\!\mathit{urves}}
\def\Polarizedstack{\mathcal{P}\!\mathit{olarized}}
\def\Complexesstack{\mathcal{C}\!\mathit{omplexes}}
% \Pic is the operator that assigns to X its picard group, usage \Pic(X)
% \Picardstack_{X/B} denotes the Picard stack of X over B
% \Picardfunctor_{X/B} denotes the Picard functor of X over B
\def\Pic{\mathop{\mathrm{Pic}}\nolimits}
\def\Picardstack{\mathcal{P}\!\mathit{ic}}
\def\Picardfunctor{\mathrm{Pic}}
\def\Deformationcategory{\mathcal{D}\!\mathit{ef}}
\usepackage{cleveref}
\crefname{originalreference}{original reference}{original references}
\Crefname{originalreference}{Original reference}{Original references}
\makeatletter
\newlabel{algebraic-lemma-characterize-representable-by-algebraic-spaces}{{algebraic-lemma-characterize-representable-by-algebraic-spaces}{source}{}{sourcepointer-1}{}}
\newlabel{algebraic-lemma-characterize-representable-by-algebraic-spaces@cref}{{[originalreference][0][]algebraic-lemma-characterize-representable-by-algebraic-spaces}{[1][0][]source}{}{}{}}
\makeatother
\crefname{originalreference}{original reference}{original references}
\Crefname{originalreference}{Original reference}{Original references}
\makeatletter
\newlabel{algebraic-lemma-map-space-into-stack}{{algebraic-lemma-map-space-into-stack}{source}{}{sourcepointer-2}{}}
\newlabel{algebraic-lemma-map-space-into-stack@cref}{{[originalreference][0][]algebraic-lemma-map-space-into-stack}{[1][0][]source}{}{}{}}
\makeatother
\crefname{originalreference}{original reference}{original references}
\Crefname{originalreference}{Original reference}{Original references}
\makeatletter
\newlabel{criteria-lemma-flat-quotient-flat-presentation}{{criteria-lemma-flat-quotient-flat-presentation}{source}{}{sourcepointer-3}{}}
\newlabel{criteria-lemma-flat-quotient-flat-presentation@cref}{{[originalreference][0][]criteria-lemma-flat-quotient-flat-presentation}{[1][0][]source}{}{}{}}
\makeatother
\crefname{originalreference}{original reference}{original references}
\Crefname{originalreference}{Original reference}{Original references}
\makeatletter
\newlabel{criteria-theorem-flat-groupoid-gives-algebraic-stack}{{criteria-theorem-flat-groupoid-gives-algebraic-stack}{source}{}{sourcepointer-4}{}}
\newlabel{criteria-theorem-flat-groupoid-gives-algebraic-stack@cref}{{[originalreference][0][]criteria-theorem-flat-groupoid-gives-algebraic-stack}{[1][0][]source}{}{}{}}
\makeatother
\crefname{originallemma}{original lemma}{original lemmas}
\Crefname{originallemma}{Original lemma}{Original lemmas}
\makeatletter
\newlabel{lemma-enough-finite-type-points}{{lemma-enough-finite-type-points}{source}{}{sourcepointer-5}{}}
\newlabel{lemma-enough-finite-type-points@cref}{{[originallemma][0][]lemma-enough-finite-type-points}{[1][0][]source}{}{}{}}
\makeatother
\crefname{originallemma}{original lemma}{original lemmas}
\Crefname{originallemma}{Original lemma}{Original lemmas}
\makeatletter
\newlabel{lemma-finite-type-points-morphism}{{lemma-finite-type-points-morphism}{source}{}{sourcepointer-6}{}}
\newlabel{lemma-finite-type-points-morphism@cref}{{[originallemma][0][]lemma-finite-type-points-morphism}{[1][0][]source}{}{}{}}
\makeatother
\crefname{originallemma}{original lemma}{original lemmas}
\Crefname{originallemma}{Original lemma}{Original lemmas}
\makeatletter
\newlabel{lemma-identify-finite-type-points}{{lemma-identify-finite-type-points}{source}{}{sourcepointer-7}{}}
\newlabel{lemma-identify-finite-type-points@cref}{{[originallemma][0][]lemma-identify-finite-type-points}{[1][0][]source}{}{}{}}
\makeatother
\crefname{originalreference}{original reference}{original references}
\Crefname{originalreference}{Original reference}{Original references}
\makeatletter
\newlabel{spaces-descent-lemma-locally-finite-presentation-fppf-local-source}{{spaces-descent-lemma-locally-finite-presentation-fppf-local-source}{source}{}{sourcepointer-8}{}}
\newlabel{spaces-descent-lemma-locally-finite-presentation-fppf-local-source@cref}{{[originalreference][0][]spaces-descent-lemma-locally-finite-presentation-fppf-local-source}{[1][0][]source}{}{}{}}
\makeatother
\crefname{originalreference}{original reference}{original references}
\Crefname{originalreference}{Original reference}{Original references}
\makeatletter
\newlabel{spaces-morphisms-lemma-composition-finite-type}{{spaces-morphisms-lemma-composition-finite-type}{source}{}{sourcepointer-9}{}}
\newlabel{spaces-morphisms-lemma-composition-finite-type@cref}{{[originalreference][0][]spaces-morphisms-lemma-composition-finite-type}{[1][0][]source}{}{}{}}
\makeatother
\crefname{originalreference}{original reference}{original references}
\Crefname{originalreference}{Original reference}{Original references}
\makeatletter
\newlabel{spaces-morphisms-lemma-noetherian-finite-type-finite-presentation}{{spaces-morphisms-lemma-noetherian-finite-type-finite-presentation}{source}{}{sourcepointer-10}{}}
\newlabel{spaces-morphisms-lemma-noetherian-finite-type-finite-presentation@cref}{{[originalreference][0][]spaces-morphisms-lemma-noetherian-finite-type-finite-presentation}{[1][0][]source}{}{}{}}
\makeatother
\crefname{originalreference}{original reference}{original references}
\Crefname{originalreference}{Original reference}{Original references}
\makeatletter
\newlabel{stacks-properties-lemma-monomorphism}{{stacks-properties-lemma-monomorphism}{source}{}{sourcepointer-11}{}}
\newlabel{stacks-properties-lemma-monomorphism@cref}{{[originalreference][0][]stacks-properties-lemma-monomorphism}{[1][0][]source}{}{}{}}
\makeatother
\crefname{originalreference}{original reference}{original references}
\Crefname{originalreference}{Original reference}{Original references}
\makeatletter
\newlabel{stacks-properties-lemma-residual-gerbe-unique}{{stacks-properties-lemma-residual-gerbe-unique}{source}{}{sourcepointer-12}{}}
\newlabel{stacks-properties-lemma-residual-gerbe-unique@cref}{{[originalreference][0][]stacks-properties-lemma-residual-gerbe-unique}{[1][0][]source}{}{}{}}
\makeatother
\crefname{originalreference}{original reference}{original references}
\Crefname{originalreference}{Original reference}{Original references}
\makeatletter
\newlabel{stacks-properties-section-properties-morphisms}{{stacks-properties-section-properties-morphisms}{source}{}{sourcepointer-13}{}}
\newlabel{stacks-properties-section-properties-morphisms@cref}{{[originalreference][0][]stacks-properties-section-properties-morphisms}{[1][0][]source}{}{}{}}
\makeatother
\title{Finite type points are exactly the points through which the residual gerbe includes as a morphism of finite type}
\author{The Stacks Project Authors}
\date{stacks-project:stacks-morphisms.tex:local-snapshot-20261012; GNU FDL 1.2 or later}
\begin{document}
\maketitle
\noindent\textit{Editorial note.} All delivered mathematical text is the author's own, copied from the pinned Stacks Project chapter source sources/stacks-algebraic-stacks-morphisms/pinned-original-source.tex; no statement or proof step was rewritten, reordered or removed.\par\medskip
\noindent\textit{Editorial note.} Preamble provenance: the Stacks Project keeps one single shared preamble for every chapter; this card takes it directly from the staged sources/stacks-algebraic-stacks-morphisms/preamble.tex as three verbatim spans (lines 8-12, 14-18 and 154-217), so no derived file was written and nothing was added or redefined. The wrapper class and the base packages are supplied by the assembler.\par\medskip
\noindent\textit{Editorial note.} Bibliography provenance: neither the counted statement nor its proof contains a citation (this lane rescans the whole span before writing the card), so the card carries no bibliography and no entry was reconstructed.\par\medskip
\noindent\textit{Editorial note.} Selection provenance: this card was opened under the file-based control contract (CONTROL.md), with the mechanical pre-card gate run on its statement and proof interval before the card was written.\par\medskip
\noindent\textit{Allowed standard prerequisites (editorial).}\begin{itemize}
\item The description of a finite type point through an affine scheme mapping smoothly to the stack with a closed point over it. Invoked at the start of the counted proof (source lines 3824-3827) as a pointer to a lemma of this chapter. Produces the presentation used to build the residual gerbe.
\item The groupoid presentation attached to a point of an algebraic stack, the construction of the quotient stack of a flat groupoid, and the uniqueness of the residual gerbe. All are invoked inside the counted proof (source lines 3828-3831, 3848-3850, 3863-3867) as cross-chapter pointers. Builds the stack whose underlying space is a singleton and identifies it with the residual gerbe.
\item Descent of the property of being locally of finite type along a surjective flat morphism of finite presentation, and the identification of finite type points through morphisms. Invoked inside the counted proof (source lines 3894-3896 and 3904-3907) as pointers to results of other chapters and of this chapter. Transfers the finiteness property to the morphism from the residual gerbe and closes the circle of implications.
\end{itemize}\medskip
\ifdefined\hypertarget\hypertarget{sourcepointer-1}{}\fi
\noindent\textit{Source reference.} Original reference algebraic-lemma-characterize-representable-by-algebraic-spaces: sources/stacks-algebraic-stacks-morphisms/pinned-original-source.tex, source lines 1720--1720. Navigation pointer only: the counted text refers to the Stacks item algebraic-lemma-characterize-representable-by-algebraic-spaces, which belongs to another chapter of the project and is not part of this card. The pointer carries the target's own label as its number because this chapter's staged snapshot does not contain the other chapter's tag table; its content is not reproduced here.\par\medskip
\ifdefined\hypertarget\hypertarget{sourcepointer-2}{}\fi
\noindent\textit{Source reference.} Original reference algebraic-lemma-map-space-into-stack: sources/stacks-algebraic-stacks-morphisms/pinned-original-source.tex, source lines 3831--3831. Navigation pointer only: the counted text refers to the Stacks item algebraic-lemma-map-space-into-stack, which belongs to another chapter of the project and is not part of this card. The pointer carries the target's own label as its number because this chapter's staged snapshot does not contain the other chapter's tag table; its content is not reproduced here.\par\medskip
\ifdefined\hypertarget\hypertarget{sourcepointer-3}{}\fi
\noindent\textit{Source reference.} Original reference criteria-lemma-flat-quotient-flat-presentation: sources/stacks-algebraic-stacks-morphisms/pinned-original-source.tex, source lines 3877--3877. Navigation pointer only: the counted text refers to the Stacks item criteria-lemma-flat-quotient-flat-presentation, which belongs to another chapter of the project and is not part of this card. The pointer carries the target's own label as its number because this chapter's staged snapshot does not contain the other chapter's tag table; its content is not reproduced here.\par\medskip
\ifdefined\hypertarget\hypertarget{sourcepointer-4}{}\fi
\noindent\textit{Source reference.} Original reference criteria-theorem-flat-groupoid-gives-algebraic-stack: sources/stacks-algebraic-stacks-morphisms/pinned-original-source.tex, source lines 3850--3850. Navigation pointer only: the counted text refers to the Stacks item criteria-theorem-flat-groupoid-gives-algebraic-stack, which belongs to another chapter of the project and is not part of this card. The pointer carries the target's own label as its number because this chapter's staged snapshot does not contain the other chapter's tag table; its content is not reproduced here.\par\medskip
\ifdefined\hypertarget\hypertarget{sourcepointer-5}{}\fi
\noindent\textit{Source reference.} Original lemma lemma-enough-finite-type-points: sources/stacks-algebraic-stacks-morphisms/pinned-original-source.tex, source lines 3907--3907. Navigation pointer only: the counted text refers to the Stacks item lemma-enough-finite-type-points, which belongs to another chapter of the project and is not part of this card. The pointer carries the target's own label as its number because this chapter's staged snapshot does not contain the other chapter's tag table; its content is not reproduced here.\par\medskip
\ifdefined\hypertarget\hypertarget{sourcepointer-6}{}\fi
\noindent\textit{Source reference.} Original lemma lemma-finite-type-points-morphism: sources/stacks-algebraic-stacks-morphisms/pinned-original-source.tex, source lines 3751--3751. Navigation pointer only: the counted text refers to the Stacks item lemma-finite-type-points-morphism, which belongs to another chapter of the project and is not part of this card. The pointer carries the target's own label as its number because this chapter's staged snapshot does not contain the other chapter's tag table; its content is not reproduced here.\par\medskip
\ifdefined\hypertarget\hypertarget{sourcepointer-7}{}\fi
\noindent\textit{Source reference.} Original lemma lemma-identify-finite-type-points: sources/stacks-algebraic-stacks-morphisms/pinned-original-source.tex, source lines 3757--3757. Navigation pointer only: the counted text refers to the Stacks item lemma-identify-finite-type-points, which belongs to another chapter of the project and is not part of this card. The pointer carries the target's own label as its number because this chapter's staged snapshot does not contain the other chapter's tag table; its content is not reproduced here.\par\medskip
\ifdefined\hypertarget\hypertarget{sourcepointer-8}{}\fi
\noindent\textit{Source reference.} Original reference spaces-descent-lemma-locally-finite-presentation-fppf-local-source: sources/stacks-algebraic-stacks-morphisms/pinned-original-source.tex, source lines 3896--3896. Navigation pointer only: the counted text refers to the Stacks item spaces-descent-lemma-locally-finite-presentation-fppf-local-source, which belongs to another chapter of the project and is not part of this card. The pointer carries the target's own label as its number because this chapter's staged snapshot does not contain the other chapter's tag table; its content is not reproduced here.\par\medskip
\ifdefined\hypertarget\hypertarget{sourcepointer-9}{}\fi
\noindent\textit{Source reference.} Original reference spaces-morphisms-lemma-composition-finite-type: sources/stacks-algebraic-stacks-morphisms/pinned-original-source.tex, source lines 3497--3497. Navigation pointer only: the counted text refers to the Stacks item spaces-morphisms-lemma-composition-finite-type, which belongs to another chapter of the project and is not part of this card. The pointer carries the target's own label as its number because this chapter's staged snapshot does not contain the other chapter's tag table; its content is not reproduced here.\par\medskip
\ifdefined\hypertarget\hypertarget{sourcepointer-10}{}\fi
\noindent\textit{Source reference.} Original reference spaces-morphisms-lemma-noetherian-finite-type-finite-presentation: sources/stacks-algebraic-stacks-morphisms/pinned-original-source.tex, source lines 3847--3847. Navigation pointer only: the counted text refers to the Stacks item spaces-morphisms-lemma-noetherian-finite-type-finite-presentation, which belongs to another chapter of the project and is not part of this card. The pointer carries the target's own label as its number because this chapter's staged snapshot does not contain the other chapter's tag table; its content is not reproduced here.\par\medskip
\ifdefined\hypertarget\hypertarget{sourcepointer-11}{}\fi
\noindent\textit{Source reference.} Original reference stacks-properties-lemma-monomorphism: sources/stacks-algebraic-stacks-morphisms/pinned-original-source.tex, source lines 1167--1167. Navigation pointer only: the counted text refers to the Stacks item stacks-properties-lemma-monomorphism, which belongs to another chapter of the project and is not part of this card. The pointer carries the target's own label as its number because this chapter's staged snapshot does not contain the other chapter's tag table; its content is not reproduced here.\par\medskip
\ifdefined\hypertarget\hypertarget{sourcepointer-12}{}\fi
\noindent\textit{Source reference.} Original reference stacks-properties-lemma-residual-gerbe-unique: sources/stacks-algebraic-stacks-morphisms/pinned-original-source.tex, source lines 3867--3867. Navigation pointer only: the counted text refers to the Stacks item stacks-properties-lemma-residual-gerbe-unique, which belongs to another chapter of the project and is not part of this card. The pointer carries the target's own label as its number because this chapter's staged snapshot does not contain the other chapter's tag table; its content is not reproduced here.\par\medskip
\ifdefined\hypertarget\hypertarget{sourcepointer-13}{}\fi
\noindent\textit{Source reference.} Original reference stacks-properties-section-properties-morphisms: sources/stacks-algebraic-stacks-morphisms/pinned-original-source.tex, source lines 248--248. Navigation pointer only: the counted text refers to the Stacks item stacks-properties-section-properties-morphisms, which belongs to another chapter of the project and is not part of this card. The pointer carries the target's own label as its number because this chapter's staged snapshot does not contain the other chapter's tag table; its content is not reproduced here.\par\medskip
\begin{lemma}
\label{lemma-point-finite-type-monomorphism}
Let $\mathcal{X}$ be an algebraic stack.
Let $x \in |\mathcal{X}|$. The following are equivalent:
\begin{enumerate}
\item $x$ is a finite type point,
\item there exists an algebraic stack $\mathcal{Z}$
whose underlying topological space $|\mathcal{Z}|$ is a singleton,
and a morphism $f : \mathcal{Z} \to \mathcal{X}$ which is
locally of finite type such that $\{x\} = |f|(|\mathcal{Z}|)$, and
\item the residual gerbe $\mathcal{Z}_x$ of $\mathcal{X}$ at $x$ exists
and the inclusion morphism $\mathcal{Z}_x \to \mathcal{X}$ is locally of
finite type.
\end{enumerate}
\end{lemma}

\begin{proof}
(All of the morphisms occurring in this paragraph are representable
by algebraic spaces, hence the conventions and results of
Properties of Stacks,
Section \ref{stacks-properties-section-properties-morphisms}
are applicable.)
Assume $x$ is a finite type point. Choose an affine scheme $U$,
a closed point $u \in U$, and a smooth morphism $\varphi : U \to \mathcal{X}$
with $\varphi(u) = x$, see
Lemma \ref{lemma-identify-finite-type-points}.
Set $u = \Spec(\kappa(u))$ as usual. Set $R = u \times_\mathcal{X} u$
so that we obtain a groupoid in algebraic spaces
$(u, R, s, t, c)$, see
Algebraic Stacks, Lemma \ref{algebraic-lemma-map-space-into-stack}.
The projection morphisms $R \to u$ are the compositions
$$
R = u \times_\mathcal{X} u \to
u \times_\mathcal{X} U \to
u \times_\mathcal{X} X = u
$$
where the first arrow is of finite type (a base change of the closed
immersion of schemes $u \to U$) and the second arrow is smooth (a base
change of the smooth morphism $U \to \mathcal{X}$). Hence
$s, t : R \to u$ are locally of finite type (as compositions, see
Morphisms of Spaces,
Lemma \ref{spaces-morphisms-lemma-composition-finite-type}).
Since $u$ is the spectrum of a field, it follows that
$s, t$ are flat and locally of finite presentation (by
Morphisms of Spaces, Lemma
\ref{spaces-morphisms-lemma-noetherian-finite-type-finite-presentation}).
We see that $\mathcal{Z} = [u/R]$ is an algebraic stack by
Criteria for Representability,
Theorem \ref{criteria-theorem-flat-groupoid-gives-algebraic-stack}.
By
Algebraic Stacks, Lemma \ref{algebraic-lemma-map-space-into-stack}
we obtain a canonical morphism
$$
f : \mathcal{Z} \longrightarrow \mathcal{X}
$$
which is fully faithful. Hence this morphism is representable by
algebraic spaces, see
Algebraic Stacks, Lemma
\ref{algebraic-lemma-characterize-representable-by-algebraic-spaces}
and a monomorphism, see
Properties of Stacks, Lemma \ref{stacks-properties-lemma-monomorphism}.
It follows that the residual gerbe $\mathcal{Z}_x \subset \mathcal{X}$
of $\mathcal{X}$ at $x$ exists and that $f$ factors through an
equivalence $\mathcal{Z} \to \mathcal{Z}_x$, see
Properties of Stacks, Lemma
\ref{stacks-properties-lemma-residual-gerbe-unique}.
By construction the diagram
$$
\xymatrix{
u \ar[d] \ar[r] & U \ar[d] \\
\mathcal{Z} \ar[r]^f & \mathcal{X}
}
$$
is commutative. By
Criteria for Representability,
Lemma \ref{criteria-lemma-flat-quotient-flat-presentation}
the left vertical arrow is surjective, flat, and locally of finite
presentation. Consider
$$
\xymatrix{
u \times_\mathcal{X} U \ar[d] \ar[r] &
\mathcal{Z} \times_\mathcal{X} U \ar[r] \ar[d] & U \ar[d] \\
u \ar[r] & \mathcal{Z} \ar[r]^f & \mathcal{X}
}
$$
As $u \to \mathcal{X}$ is locally of finite type, we see that the base change
$u \times_\mathcal{X} U \to U$ is locally of finite type. Moreover,
$u \times_\mathcal{X} U \to \mathcal{Z} \times_\mathcal{X} U$ is
surjective, flat, and locally of finite presentation as a base change of
$u \to \mathcal{Z}$. Thus
$\{u \times_\mathcal{X} U \to \mathcal{Z} \times_\mathcal{X} U\}$
is an fppf covering of algebraic spaces, and we conclude that
$\mathcal{Z} \times_\mathcal{X} U \to U$ is locally of finite type by
Descent on Spaces, Lemma
\ref{spaces-descent-lemma-locally-finite-presentation-fppf-local-source}.
By definition this means that $f$ is locally of finite type (because the
vertical arrow $\mathcal{Z} \times_\mathcal{X} U \to \mathcal{Z}$ is smooth
as a base change of $U \to \mathcal{X}$ and surjective as $\mathcal{Z}$ has
only one point). Since $\mathcal{Z} = \mathcal{Z}_x$ we see that (3) holds.

\medskip\noindent
It is clear that (3) implies (2).
If (2) holds then $x$ is a finite type point of $\mathcal{X}$ by
Lemma \ref{lemma-finite-type-points-morphism}
and
Lemma \ref{lemma-enough-finite-type-points}
to see that $\mathcal{Z}_{\text{ft-pts}}$ is nonempty, i.e., the
unique point of $\mathcal{Z}$ is a finite type point of $\mathcal{Z}$.
\end{proof}

\end{document}
